El presente informe corresponde al quinto y último \textit{checkpoint} del proyecto integrador, cuyo objetivo fue completar el diseño, la implementación y la validación experimental de una fuente conmutada tipo \textit{buck} operando a lazo cerrado. A diferencia de las entregas anteriores, centradas principalmente en el dimensionamiento teórico y en la verificación parcial de bloques individuales, en esta etapa se integraron todos los subsistemas necesarios para obtener un prototipo funcional.

En primer lugar, se retomó la caracterización del inductor construido, ya que este componente resulta determinante para el comportamiento dinámico del convertidor. A partir de las mediciones realizadas se verificó su inductancia efectiva, su margen frente a la saturación y su compatibilidad con las condiciones de operación previstas para la fuente.

Luego se desarrolló la etapa de potencia y control del convertidor. Esto incluyó la modificación del generador PWM, la selección de los diodos de rueda libre, la incorporación de una carga fantasma, el análisis de los capacitores de entrada y salida, el diseño del \textit{gate driver} y la implementación de una red de compensación para estabilizar el sistema a lazo cerrado. También se presenta el circuito final adoptado y las decisiones de diseño que permitieron adaptar el modelo teórico a una implementación práctica.

Posteriormente, se realizaron simulaciones para validar la respuesta temporal y frecuencial del sistema compensado. Estas permitieron comparar el comportamiento del convertidor antes y después de la compensación, evaluar la estabilidad del lazo y estimar el desempeño esperado frente a variaciones de carga y tensión de entrada.

Finalmente, se diseñó e implementó el PCB del prototipo y se llevaron a cabo mediciones experimentales sobre la fuente construida. En particular, se analizaron las señales del \textit{gate driver}, la conmutación de la etapa de potencia, la eficiencia, la regulación de línea y carga, y el \textit{ripple} de salida. De esta manera, el informe busca contrastar los resultados teóricos y simulados con el comportamiento real del circuito, identificando las principales diferencias observadas y las correcciones realizadas durante la puesta en marcha.